%---------------------------------------------------------------------
%
%                        frontmatter/abstract.tex
%
%---------------------------------------------------------------------
%
% The English abstract and its keywords (V.2b). Half a page, and it has to
% include the English title.
%
% The body of this report is written in English, so, unlike in the
% template, this file needs no otherlanguage wrapper.
%
%---------------------------------------------------------------------

\chapter*{Abstract}
\addcontentsline{toc}{chapter}{Abstract}

\section*{\tituloPortadaEngVal}

\todo{Write the half-page abstract once the results are in. It should
say, in this order: (1) that musical similarity has no operational
definition and that the literature substitutes genre for it; (2) the
general objective, studying similarity-learning techniques for music
analysis and graph construction, with siamese networks over three kinds
of input (spectrograms, audio embeddings, acoustic descriptors) as the
principal focus, compared against classical classification and
similarity techniques and against a fuzzy block for gradual relationships
(O1--O4); (3) that the graphs are evaluated at a comparative level
(do representations differ significantly?) and an absolute level (is
each graph good on its own, and does it agree with human judgements?),
with FMA and MagnaTagATune, using Recall@K, MAP and nDCG (O5); (4) the
headline result, including how far the rankings agree with human
similarity judgements and whether the fuzzy block helps; and (5) the
conclusion, answering the main research question stated in Chapter
\ref{cap:introduction}.}

\section*{Keywords}

\noindent \todo{At most 10 keywords, separated by commas. Proposal: music
information retrieval, music similarity, siamese networks, similarity
graphs, metric learning, classical classification, fuzzy logic,
evaluation, human similarity judgements, MERT.}
